\Askhsh{1773}{4}
\begin{ylh}{1}
    \item 5.2. Παραλληλόγραμμα
    \item 5.4. Ρόμβος
    \item 5.6. Εφαρμογές στα τρίγωνα.
\end{ylh}
Δίνεται τετράπλευρο $\mathrm{A}\mathrm{B}\Gamma\Delta$ με $\mathrm{A}\Delta = \mathrm{B}\Gamma$. Αν $\mathrm{E}, \Lambda, \mathrm{Z}, \mathrm{K}, \mathrm{N}, \mathrm{M}$ είναι τα μέσα των $\mathrm{A}\mathrm{B}, \mathrm{B}\Gamma, \Gamma\Delta, \Delta\mathrm{A}, \Delta\mathrm{B}$ και $\mathrm{A}\Gamma$ αντίστοιχα, να αποδείξετε ότι:
\begin{alist}
    \item Το τετράπλευρο $\mathrm{E}\mathrm{M}\mathrm{Z}\mathrm{N}$ ρόμβος. \monades{8}
    \item Η $\mathrm{E}\mathrm{Z}$ είναι μεσοκάθετος του ευθύγραμμου τμήματος $\mathrm{M}\mathrm{N}$. \monades{7}
    \item $\mathrm{K}\mathrm{E} = \mathrm{Z}\Lambda$. \monades{5}
    \item Τα ευθύγραμμα τμήματα $\mathrm{K}\Lambda, \mathrm{M}\mathrm{N}, \mathrm{E}\mathrm{Z}$ διέρχονται από το ίδιο σημείο. \monades{5}
\end{alist}
